\aufzaehlungvier{Hal ini mengikuti dari deskripsi eksplisit turunan vertikal suatu koneksi linear; lihat
Catatan.
}{Jelas dari (1).
}{Hal ini mengikuti dari (2) dan
Lema.
}{Kita tuliskan

\mavergleichskette
{\vergleichskette
{ V
}
{ = }{ \sum_{i = 1}^n {h_i } \partial_i
}
{  }{
}
{    }{
}
{   }{
}
}
{}{}{}
dan

\mavergleichskette
{\vergleichskette
{ W
}
{ = }{ \sum_{j = 1}^n {f_j } \partial_j
}
{  }{
}
{    }{
}
{   }{
}
}
{}{}{.}
Menurut
Lema
dan bagian-bagian yang telah dibuktikan,

\mavergleichskettealign
{\vergleichskettealign
{ \nabla_V W
}
{ =} { \nabla_{\sum_{i = 1}^n {h_i } \partial_i }   \left(  \sum_{j = 1}^n {f_j } \partial_j   \right)
}
{ =} { \sum_{j = 1}^n \left(  \sum_{i = 1}^n h_i  \partial_i f_j   \right) \partial_j
}
{ } {
}
{ } {
}
}
{}
{}{.}
Pada suatu titik

\mavergleichskette
{\vergleichskette
{ P
}
{ \in }{ U
}
{  }{
}
{    }{
}
{   }{
}
}
{}{}{}
nilai ini adalah

\mathdisp {\sum_{j  = 1}^n \left(  \sum_{i = 1}^n h_i (P) ( \partial_i f_j)(P)   \right) \partial_j} { . }

Di sisi lain,

\mavergleichskettealign
{\vergleichskettealign
{ \left(  D W   \right)_{ P } \left(   V(P)  \right)
}
{ =} { \left(  (\partial_j f_i)(P)  \right)_{ji} \begin{pmatrix} h_1(P) \\ \vdots\\  h_n(P)   \end{pmatrix}
}
{ =} { \begin{pmatrix}  \sum_{i = 1}^n h_i (P) \left(  \partial_i f_1  \right)(P)  \\ \vdots\\  \sum_{i = 1}^n h_i (P)  \left(  \partial_i f_n  \right) (P)   \end{pmatrix}
}
{ } {
}
{ } {
}
}
{}
{}{.}
}
